% w12-tilt-geometry.tex
%   12주차 §12-1. 추진기 넷의 자리와 틸팅 각 delta 의 정의, 그리고 틸팅 추진기
%   하나가 B 에 **열 둘**을 내놓는 이유.
%
%   축척 / scale
%       1 m = 36 mm. 화면 기저는 평면도 — 화면 위가 선체 x_b, 화면 오른쪽이 y_b.
%       따라서 선체좌표 (x, y) 는 화면 (X, Y) = (36 y, 36 x) 에 놓인다.
%
%   회전 방향 / the sense of rotation            (CLAUDE.md 4-6-2)
%       z_b 가 아래이므로 +N 은 선수를 우현으로 돌리고, 위에서 내려다본 이 그림
%       (x_b 위, y_b 오른쪽)에서 그것은 **시계방향**이다. delta 도 같다 — x_b 에서
%       y_b 로 재므로 화면에서는 "위에서 오른쪽으로", 곧 시계방향이다.
%       TikZ 의 arc 는 끝 각이 **작을 때** 시계이므로 둘 다 arc(90:55:..) ·
%       arc(330:210:..) 처럼 **내려가는** 각으로 적었다.
%
%   l_f 가 선체보다 긴 것 / the forward arm is longer than the hull
%       otter_config.m 의 l_f = 1.5 m 는 **모멘트 팔**이고 선체는 L = 2.0 m
%       (x in [-1, +1]) 다. 6주차가 같은 상수를 공유하므로 한쪽만 고칠 수 없다.
%       사용자 지시(2026-10-13)는 **유닛을 선체 밖에 그리지 말라**는 것이었다 —
%       그래서 6주차 w06-alloc-geometry.tex 와 **같은 처리**를 한다.
%       유닛 넷은 폰툰의 실제 끝(x = +-1.0 m)에 그리고, l_f 는 치수선으로만
%       선수 앞 0.5 m 까지 뻗어 "팔"이라고 밝힌다. 선미는 l_a = 1.0 m 가 실제
%       자리와 같으므로 어긋나는 것은 선수 하나뿐이다.
%       6주차와 같은 문구: "Read the arms, not the hull stations."

\documentclass[border=5pt]{standalone}
% gnc-style.tex — 이 볼트의 TikZ 그림이 공유하는 스타일.
%
% 저널 그림 규약을 스타일 하나로 굳혀 둔다. 그림마다 선 굵기나 화살촉을
% 다시 정하지 않는다 — 그것이 그림이 제각각으로 보이는 원인이다.
%
% 쓰는 법:  % gnc-style.tex — 이 볼트의 TikZ 그림이 공유하는 스타일.
%
% 저널 그림 규약을 스타일 하나로 굳혀 둔다. 그림마다 선 굵기나 화살촉을
% 다시 정하지 않는다 — 그것이 그림이 제각각으로 보이는 원인이다.
%
% 쓰는 법:  % gnc-style.tex — 이 볼트의 TikZ 그림이 공유하는 스타일.
%
% 저널 그림 규약을 스타일 하나로 굳혀 둔다. 그림마다 선 굵기나 화살촉을
% 다시 정하지 않는다 — 그것이 그림이 제각각으로 보이는 원인이다.
%
% 쓰는 법:  \input{gnc-style}  (figures/src/ 안에서)

\usepackage{tikz}
\usepackage{amsmath}
\usepackage{xcolor}
\usetikzlibrary{arrows.meta, positioning, calc, fit, backgrounds, decorations.pathmorphing}

% ---- 색 --------------------------------------------------------------------
% 색은 강조 하나에만 쓴다. 나머지는 검정.
\definecolor{ink}{HTML}{1A1A1A}
\definecolor{accent}{HTML}{7C3AED}
\definecolor{warn}{HTML}{B91C1C}
\definecolor{soft}{HTML}{666666}

% 기하 그림(좌표계·유도)에서만 쓰는 넷. 블록선도는 위의 흑백 + accent 로 충분하다.
\definecolor{pathc}{HTML}{0D9488}   % 계획 경로 · 항적
\definecolor{errc}{HTML}{C2410C}    % 오차 (cross-track 등)
\definecolor{alng}{HTML}{B45309}    % along-track
\definecolor{flow}{HTML}{0369A1}    % 조류 · 유속

% ---- 선과 화살촉 -----------------------------------------------------------
\tikzset{
  % 신호선. 화살촉은 작고 뾰족한 것 하나뿐이다.
  sig/.style   = {ink, line width=0.5pt, -{Stealth[length=4pt, width=3pt]}},
  wire/.style  = {ink, line width=0.5pt},
  % 강조 경로 — 파선으로 구분한다. 흑백 인쇄에서도 살아야 하기 때문이다.
  asig/.style  = {accent, line width=0.5pt, densely dashed,
                  -{Stealth[length=4pt, width=3pt]}},
  awire/.style = {accent, line width=0.5pt, densely dashed},
  % 블록. 흰 바탕, 직각, 검은 테두리.
  blk/.style   = {draw=ink, line width=0.55pt, fill=white,
                  minimum width=17mm, minimum height=8mm, inner sep=2pt, font=\small},
  % 합산점
  sum/.style   = {draw=ink, line width=0.55pt, fill=white, circle,
                  inner sep=0pt, minimum size=4.6mm, font=\scriptsize},
  asum/.style  = {draw=accent, line width=0.55pt, fill=white, circle,
                  inner sep=0pt, minimum size=4.6mm, font=\scriptsize},
  % 분기점
  dot/.style   = {ink, fill=ink, circle, inner sep=0pt, minimum size=1.5mm},
  % 신호 이름 — 선 위에 작게
  sname/.style = {font=\small, inner sep=1.5pt},
  % 그림 안의 짧은 설명. 그림 안의 글은 최소로 둔다
  note/.style  = {font=\scriptsize, soft, inner sep=1.5pt},
  anote/.style = {font=\scriptsize, accent, inner sep=1.5pt},
  % 축
  axis/.style  = {ink, line width=0.5pt},
  tick/.style  = {font=\scriptsize, soft},
}

% 캡션 한 줄 — 그림 아래 가는 선과 함께
\newcommand{\gnccaption}[2]{%
  \node[anchor=north west, text width=#1, align=left, font=\scriptsize, ink]
        at (current bounding box.south west) {\vspace{2pt}#2};}
  (figures/src/ 안에서)

\usepackage{tikz}
\usepackage{amsmath}
\usepackage{xcolor}
\usetikzlibrary{arrows.meta, positioning, calc, fit, backgrounds, decorations.pathmorphing}

% ---- 색 --------------------------------------------------------------------
% 색은 강조 하나에만 쓴다. 나머지는 검정.
\definecolor{ink}{HTML}{1A1A1A}
\definecolor{accent}{HTML}{7C3AED}
\definecolor{warn}{HTML}{B91C1C}
\definecolor{soft}{HTML}{666666}

% 기하 그림(좌표계·유도)에서만 쓰는 넷. 블록선도는 위의 흑백 + accent 로 충분하다.
\definecolor{pathc}{HTML}{0D9488}   % 계획 경로 · 항적
\definecolor{errc}{HTML}{C2410C}    % 오차 (cross-track 등)
\definecolor{alng}{HTML}{B45309}    % along-track
\definecolor{flow}{HTML}{0369A1}    % 조류 · 유속

% ---- 선과 화살촉 -----------------------------------------------------------
\tikzset{
  % 신호선. 화살촉은 작고 뾰족한 것 하나뿐이다.
  sig/.style   = {ink, line width=0.5pt, -{Stealth[length=4pt, width=3pt]}},
  wire/.style  = {ink, line width=0.5pt},
  % 강조 경로 — 파선으로 구분한다. 흑백 인쇄에서도 살아야 하기 때문이다.
  asig/.style  = {accent, line width=0.5pt, densely dashed,
                  -{Stealth[length=4pt, width=3pt]}},
  awire/.style = {accent, line width=0.5pt, densely dashed},
  % 블록. 흰 바탕, 직각, 검은 테두리.
  blk/.style   = {draw=ink, line width=0.55pt, fill=white,
                  minimum width=17mm, minimum height=8mm, inner sep=2pt, font=\small},
  % 합산점
  sum/.style   = {draw=ink, line width=0.55pt, fill=white, circle,
                  inner sep=0pt, minimum size=4.6mm, font=\scriptsize},
  asum/.style  = {draw=accent, line width=0.55pt, fill=white, circle,
                  inner sep=0pt, minimum size=4.6mm, font=\scriptsize},
  % 분기점
  dot/.style   = {ink, fill=ink, circle, inner sep=0pt, minimum size=1.5mm},
  % 신호 이름 — 선 위에 작게
  sname/.style = {font=\small, inner sep=1.5pt},
  % 그림 안의 짧은 설명. 그림 안의 글은 최소로 둔다
  note/.style  = {font=\scriptsize, soft, inner sep=1.5pt},
  anote/.style = {font=\scriptsize, accent, inner sep=1.5pt},
  % 축
  axis/.style  = {ink, line width=0.5pt},
  tick/.style  = {font=\scriptsize, soft},
}

% 캡션 한 줄 — 그림 아래 가는 선과 함께
\newcommand{\gnccaption}[2]{%
  \node[anchor=north west, text width=#1, align=left, font=\scriptsize, ink]
        at (current bounding box.south west) {\vspace{2pt}#2};}
  (figures/src/ 안에서)

\usepackage{tikz}
\usepackage{amsmath}
\usepackage{xcolor}
\usetikzlibrary{arrows.meta, positioning, calc, fit, backgrounds, decorations.pathmorphing}

% ---- 색 --------------------------------------------------------------------
% 색은 강조 하나에만 쓴다. 나머지는 검정.
\definecolor{ink}{HTML}{1A1A1A}
\definecolor{accent}{HTML}{7C3AED}
\definecolor{warn}{HTML}{B91C1C}
\definecolor{soft}{HTML}{666666}

% 기하 그림(좌표계·유도)에서만 쓰는 넷. 블록선도는 위의 흑백 + accent 로 충분하다.
\definecolor{pathc}{HTML}{0D9488}   % 계획 경로 · 항적
\definecolor{errc}{HTML}{C2410C}    % 오차 (cross-track 등)
\definecolor{alng}{HTML}{B45309}    % along-track
\definecolor{flow}{HTML}{0369A1}    % 조류 · 유속

% ---- 선과 화살촉 -----------------------------------------------------------
\tikzset{
  % 신호선. 화살촉은 작고 뾰족한 것 하나뿐이다.
  sig/.style   = {ink, line width=0.5pt, -{Stealth[length=4pt, width=3pt]}},
  wire/.style  = {ink, line width=0.5pt},
  % 강조 경로 — 파선으로 구분한다. 흑백 인쇄에서도 살아야 하기 때문이다.
  asig/.style  = {accent, line width=0.5pt, densely dashed,
                  -{Stealth[length=4pt, width=3pt]}},
  awire/.style = {accent, line width=0.5pt, densely dashed},
  % 블록. 흰 바탕, 직각, 검은 테두리.
  blk/.style   = {draw=ink, line width=0.55pt, fill=white,
                  minimum width=17mm, minimum height=8mm, inner sep=2pt, font=\small},
  % 합산점
  sum/.style   = {draw=ink, line width=0.55pt, fill=white, circle,
                  inner sep=0pt, minimum size=4.6mm, font=\scriptsize},
  asum/.style  = {draw=accent, line width=0.55pt, fill=white, circle,
                  inner sep=0pt, minimum size=4.6mm, font=\scriptsize},
  % 분기점
  dot/.style   = {ink, fill=ink, circle, inner sep=0pt, minimum size=1.5mm},
  % 신호 이름 — 선 위에 작게
  sname/.style = {font=\small, inner sep=1.5pt},
  % 그림 안의 짧은 설명. 그림 안의 글은 최소로 둔다
  note/.style  = {font=\scriptsize, soft, inner sep=1.5pt},
  anote/.style = {font=\scriptsize, accent, inner sep=1.5pt},
  % 축
  axis/.style  = {ink, line width=0.5pt},
  tick/.style  = {font=\scriptsize, soft},
}

% 캡션 한 줄 — 그림 아래 가는 선과 함께
\newcommand{\gnccaption}[2]{%
  \node[anchor=north west, text width=#1, align=left, font=\scriptsize, ink]
        at (current bounding box.south west) {\vspace{2pt}#2};}


\begin{document}
\begin{tikzpicture}[x=1mm, y=1mm, font=\scriptsize]

% ---- 축척 상수 / the scale ----------------------------------------------
\def\YP{14.22}      % y_pont = 0.395 m
\def\HL{36}         % 선체 반길이 1.0 m — 유닛이 앉는 실제 자리
\def\ARM{54}        % 선수 모멘트 팔 l_f = 1.50 m 의 끝 (치수선만)

% =======================================================================
% 1. 선체 — 폰툰 둘 / the hull: two pontoons
% =======================================================================
\fill[ink, opacity=0.10, rounded corners=3pt] (-18.72,-36) rectangle (-9.72,36);
\fill[ink, opacity=0.10, rounded corners=3pt] ( 9.72,-36) rectangle (18.72,36);
\draw[ink, line width=0.5pt, rounded corners=3pt] (-18.72,-36) rectangle (-9.72,36);
\draw[ink, line width=0.5pt, rounded corners=3pt] ( 9.72,-36) rectangle (18.72,36);
\draw[soft, line width=0.4pt] (-9.72, 36) -- (9.72, 36);
\draw[soft, line width=0.4pt] (-9.72,-36) -- (9.72,-36);
\node[soft, anchor=north, font=\scriptsize] at (0, 34.5) {bow};

% =======================================================================
% 2. 선체좌표계 / the body frame
% =======================================================================
\draw[ink, line width=0.5pt, -{Stealth[length=4pt,width=3pt]}]
      (0,0) -- (0,26) node[anchor=south west, inner sep=1pt] {$x_b$};
\draw[ink, line width=0.5pt, -{Stealth[length=4pt,width=3pt]}]
      (0,0) -- (26,0) node[anchor=north west, inner sep=1pt] {$y_b$};
\fill[ink] (0,0) circle (0.8);
\node[ink, anchor=north east, inner sep=1.5pt] at (0,0) {$o_b$};
\draw[soft, dotted, line width=0.4pt] (0,0) -- (0,-48);

% +N 은 시계방향 / a positive yaw moment is clockwise on this view
%   arc(330:210:7) — 끝 각이 작으므로 시계 / decreasing angle, hence clockwise
\draw[warn, line width=0.6pt, -{Stealth[length=4pt,width=3pt]}]
      (6.06,-3.5) arc (330:210:7);
\node[warn, anchor=north, inner sep=1.5pt] at (0,-12) {$\tau_N > 0$};

% =======================================================================
% 3. 추진기 넷 / the four units — 폰툰의 실제 끝에 앉는다
% =======================================================================
\foreach \p/\q/\lab in {-\YP/\HL/1, \YP/\HL/2, -\YP/-\HL/3, \YP/-\HL/4}{
  \node[circle, draw=ink, fill=white, line width=0.5pt,
        inner sep=0pt, minimum size=6mm] at (\p,\q) {\lab};
}

% 추력 — 넷 모두 같은 delta = 35 deg / all four at the same delta
%   방향 = (sin d, cos d) * 20 mm = (11.47, 16.38)
\foreach \p/\q in {-\YP/\HL, \YP/\HL, -\YP/-\HL, \YP/-\HL}{
  \draw[accent, line width=0.8pt, -{Stealth[length=5pt,width=4pt]}]
        (\p,\q) -- ++(11.47,16.38);
}

% 1번만 각을 표시한다 / only unit 1 carries the angle annotation
\draw[soft, dashed, line width=0.4pt] (-\YP,\HL) -- (-\YP,60)
      node[anchor=south west, inner sep=1.5pt, soft] {$\parallel x_b$};
\draw[accent, line width=0.5pt] (-14.22,51) arc (90:55:15);
\node[accent, inner sep=0pt] at (-11.2,45.5) {$\delta_1$};
\node[accent, anchor=south west, inner sep=1.5pt] at (-2.0,52.4)
      {$T_1\,\mathbf{e}_1$};

% =======================================================================
% 4. 모멘트 팔 / the moment arms
% =======================================================================
%  선수 팔만 선체보다 길다 — 치수선으로만 뻗고 유닛은 선체 위에 있다
\draw[soft, dotted, line width=0.4pt] (\YP,39) -- (\YP,\ARM);
\draw[soft, dotted, line width=0.4pt] (\YP,\ARM) -- (32,\ARM);
\draw[soft, dotted, line width=0.4pt] ( 1.0,  0 ) -- (32,  0 );
\draw[soft, dotted, line width=0.4pt] (18.0,-36) -- (32,-36);

\draw[accent, line width=0.6pt, |-|] (32,0) -- (32,\ARM);
\node[accent, anchor=west, inner sep=1.5pt, text width=26mm] at (33.5,27)
      {$l_f = 1.50$ m\\ moment arm};
\draw[accent, line width=0.6pt, |-|] (32,-36) -- (32,0);
\node[accent, anchor=west, inner sep=1.5pt, text width=26mm] at (33.5,-18)
      {$l_a = 1.00$ m\\ at the transom};

\draw[soft, dotted, line width=0.4pt] (\YP,-39) -- (\YP,-48);
\draw[accent, line width=0.6pt, |-|] (0,-48) -- (\YP,-48);
\node[accent, anchor=north, inner sep=1.5pt] at (7.1,-50)
      {$y_{\mathrm{pont}} = 0.395$ m};

% =======================================================================
% 5. 오른쪽 설명 / the panel
% =======================================================================
\begin{scope}[shift={(66,0)}]

\node[ink, anchor=north west, text width=84mm, align=left, font=\scriptsize]
      at (0,58) {%
  \textbf{Two columns, not one}\\[2pt]
  A fixed thruster spends its thrust $T_i$ along a direction the hull
  fixes. A tilting unit chooses the direction as well, so the pair
  $(T_i, \delta_i)$ carries two independent numbers:
  \\[3pt]
  \quad $\mathbf{e}_i = \begin{bmatrix}\cos\delta_i\\ \sin\delta_i\end{bmatrix}$,
  \qquad
  $\mathbf{f}_i = T_i\,\mathbf{e}_i
     = \begin{bmatrix}f_{ix}\\ f_{iy}\end{bmatrix}$
  \\[4pt]
  The column rule of Week 3 takes a force $\mathbf{f}$ applied at
  $(x_i, y_i)$ and returns what it does to the hull:
  \\[3pt]
  \quad $\begin{bmatrix} f_x \\ f_y \\ x_i f_y - y_i f_x \end{bmatrix}$
  \\[4pt]
  One free component, one column. A tilting unit has two, so it owns
  \emph{two} columns of $B$, and the unknown of the allocation problem is
  \\[3pt]
  \quad $\mathbf{f} = \begin{bmatrix} f_{1x} & f_{1y} & \cdots &
        f_{4x} & f_{4y}\end{bmatrix}^{\!\top}$,
  \quad eight numbers,
  \\[3pt]
  against which $\boldsymbol{\tau} = B\mathbf{f}$ is three. Stage~1 never
  solves that system: the eight numbers are given by hand, and the only
  question is whether $\boldsymbol{\tau}$ comes out as the rule says.
  \\[4pt]
  \textbf{Sense.} $\delta_i$ is measured from $x_b$ towards $y_b$. With
  $z_b$ pointing down, both a positive $\delta$ and a positive $\tau_N$
  appear \emph{clockwise} on this plan view.
  \\[4pt]
  \textbf{Read the arms, not the hull stations.} The four units are drawn
  where they sit on the pontoons, at $x_b = \pm 1.0$~m. The aft arm
  $l_a$ is that station; the forward arm $l_f = 1.50$~m reaches 0.5~m
  past the bow of the $L = 2.0$~m hull, and it is the arm, not a place.
  Week~6 carries the same constant and the same reading.
};

\end{scope}

\end{tikzpicture}
\end{document}
